\section{Related Work}

\subsection{Multimodal Misinformation Detection}

NewsCLIPpings introduced a large benchmark of authentic news images paired with deliberately mismatched captions \cite{luo2021newsclippings}. VERITE later addressed unimodal bias and topic imbalance in real-world evaluation \cite{papadopoulos2024verite}. KAMP uses single-modality, cross-modal, and knowledge-content objectives during multimodal pretraining \cite{zhang2025kamp}. Other work uses synthetic examples to improve transfer to real misinformation cases \cite{zeng2024synthetic}. Event-Radar represents visual manipulation, textual emotion, and event-level inconsistency in an event graph \cite{ma2024eventradar}. These methods provide the main no-retrieval and event-aware baselines in this study.

Uncertainty-aware fusion is another relevant content-only line. MDF-FND represents modality uncertainty and changes the contribution of text and images dynamically \cite{lv2025mdffnd}. Propagation-based methods instead use how a post spreads; RPPG-Fake generates early heterogeneous propagation paths \cite{zhang2025rppg}. Neither setting establishes the origin or historical availability of external web evidence, but both motivate controlled uncertainty and early-information comparisons.

\subsection{Evidence-Grounded Fact-Checking}

MOCHEG treats multimodal fact-checking as evidence retrieval followed by verification and explanation \cite{yao2023mocheg}. MetaSumPerceiver summarizes evidence from several multimodal documents \cite{chen2024metasum}. RED-DOT addresses noisy web evidence through relevance-aware reranking \cite{papadopoulos2025reddot}. RAGAR and MultiReflect combine multimodal language models with iterative retrieval \cite{khaliq2024ragar,kangur2025multireflect}, while SNIFFER joins internal image--text consistency with external image context \cite{qi2024sniffer}. These studies motivate strong document-level retrieval and reasoning baselines. They do not test whether a large result set reflects several independent sources or copies of one source.

\subsection{Image Context, Time, and Evidence Dependence}

The 5Pils system recovers an image's original context from web evidence and excludes evidence that follows the fact-check \cite{tonglet2024image}; COVE predicts that context before assessing a caption \cite{tonglet2025cove}. Chronological distribution shift can substantially reduce multimodal misinformation performance \cite{stepanova2023temporal}. ChronoFact handles temporal relations within complex textual claims through timeline reasoning \cite{barik2025chronofact}. These works concern context, event time, or future generalization. The present protocol focuses on a different temporal variable: the earliest independently observed availability of each evidence item relative to a claim-specific cutoff.

Source information also affects decisions under conflicting evidence \cite{ge2025conflicting}. Because source filtering can discard useful disagreement, \systemname{} separates dependence from credibility. Exact and near-duplicate content, shared imagery, direct links, and publication order support dependency clustering; publisher reputation does not determine primary-model inclusion. The joint object of study is therefore decision-time-valid independent corroboration, rather than a new graph encoder or a general source-trust score.

